\newgeometry{left=1cm,right=1cm,top=1cm,bottom=1cm,}
\begin{titlepage}
\vspace*{4cm}
\begin{center}
	\vspace{1em}\par \rule{15cm}{.1pt}
	{\HUGE Thermodynamique\\ de l’ingénieur}
	\vspace{1em}\par \rule{15cm}{.1pt}

	\vspace{4cm}
	\includegraphics[width=8cm]{images/thermodynamics.png}
	\vspace{2cm}

	{\LARGE Olivier Cleynen\\
	\ccLogo\ \ccAttribution\ \ccShareAlike}\\
	\includegraphics[width=60mm]{images/Logo_framabook_grand.png}
\end{center}

\clearpage

\vspace*{5cm}
{\parshape 1 .15\hsize.65\hsize
	Framasoft a été créé en novembre 2001 par Alexis Kauffmann. En janvier 2004 une association éponyme a vu le jour pour soutenir le développement du réseau. Pour plus d’information sur Framasoft, consulter http://www.framasoft.net.

}

{\parshape 1 .15\hsize.65\hsize
	Se démarquant de l’édition classique, les Framabooks sont dits «~livres libres~» parce qu’ils sont placés sous une licence qui permet au lecteur de disposer des mêmes libertés qu’un utilisateur de logiciels libres. Les Framabooks s’inscrivent dans cette culture des biens communs qui, à l’instar de Wikipédia, favorise la création, le partage, la diffusion et l’appropriation collective de la connaissance.

}

{\parshape 1 .15\hsize.65\hsize
	Le projet Framabook est coordonné par Christophe Masutti. Pour plus d’information, consultez http://framabook.org.

}

\vspace*{\stretch{4}}
\begin{center}
		\noindent Copyright 2014 Olivier Cleynen, Framasoft (collection Framabook)\\
		\textit{Thermodynamique de l’ingénieur} est placé sous~: Attribution -- Partage dans les Mêmes Conditions 4.0

		\noindent ISBN~: xxxxxxxxxxxxxxxxx \\
		Prix~: xx euros \\
		Dépôt légal~: 2014, \textsc{imprimeur}
\end{center}


\end{titlepage}
\restoregeometry
